\documentclass{beamer}

\usetheme{default}

\usepackage{times}

\setbeamertemplate{navigation symbols}{}
\setbeamertemplate{footline}{
	\vbox{
		\hfill
		\usebeamerfont{framenumber}\insertframenumber/\inserttotalframenumber
		\hspace{2ex}
	}
	\noindent\includegraphics[width=\textwidth]{LatexTemplates/bar.pdf}
}

\newcommand{\topBar}{
	\setbeamertemplate{headline}{
		\includegraphics[width=\textwidth]{LatexTemplates/top.pdf}
	}
	\setbeamertemplate{footline}{
	}
}

\newcommand{\sponsorBar}{
	\setbeamertemplate{footline}{
		\includegraphics[width=\textwidth]{LatexTemplates/newsponsors.pdf}
	}
}

\newcommand{\workshop}{SET WORKSHOP}
\renewcommand{\workshop}{NTU College of Public Health}
\renewcommand{\workshop}{MMED}
\renewcommand{\workshop}{DAIDD}

%% \newcommand{\figshare}{https://figshare.com/collections/International_Clinics_on_Infectious_Disease_Dynamics_and_Data/3788224}
\newcommand{\figshare}{http://www.ici3d.org/\workshop}
\newcommand{\thisyear}{2025} 

%% \newcommand{\jdauth}{Jonathan Dushoff, McMaster University}
\newcommand{\jddate}{\workshop\ \thisyear} 

\input{talkdir/bulletspace.tex}

\usepackage{talkdir/slidecomm,talkdir/commands}
\usepackage{graphicx,xspace}
\usepackage{grffile}
\usepackage{hyperref}
\usepackage{url}

\input{talkdir/section.tex}
\input{talkdir/lengths.tex}

\begin{document}



\frametitle{}

\newcommand{\jdtitle}{Heterogeneity, contact patterns and modeling options}
\newcommand{\jdsub}{}
\newcommand{\years}{2013--2020}

 \date{\jddate} \title{\jdtitle} \subtitle{\jdsub} {\topBar\maketitle}
\begin{frame}


\frametitle{Goals}\begin{itemize}[<+-| alert@+>]

\item Explain the importance of heterogeneity on patterns of disease spread\begin{itemize}[<+-| alert@+>]
\item Focus on different types of human heterogeneity\end{itemize}

\item Discuss ways in which homogeneous models fail to match observed
	dynamics

\item Use simple models to explore qualitative effects of heterogeneity on
	modeling conclusions

\item Briefly introduce some methods that are used to incorporate heterogeneity in
	models\end{itemize}
\uncover<+->\mbox{}\end{frame}

\begin{frame}


\frametitle{The resilience of infectious disease}

\begin{columns}[c] \column{0.54\textwidth} \small

1967: It's time to close the book on infectious diseases

\hfill\includegraphics[width=0.6\textwidth]{my_images/smallpox.jpg}\hfill\mbox{}

\column{0.4\textwidth} \small

\hfill\includegraphics[width=0.6\textwidth]{my_images/headlines.jpg}\hfill\mbox{}

\hfill\includegraphics[width=0.6\textwidth]{my_images/book.jpg}\hfill\mbox{}

\end{columns}
\uncover<+->\mbox{}\end{frame}

\begin{frame}


\frametitle{Pathogen evolution}

\hfill\includegraphics[width=0.4\textwidth]{my_images/influenza.jpg}\hfill\mbox{}

\hfill\includegraphics[width=0.4\textwidth]{my_images/tb.jpg}\hfill\mbox{}
\uncover<+->\mbox{}\end{frame}

\begin{frame}


\frametitle{Human heterogeneity}

\hfill\includegraphics[height=0.42\textheight]{my_images/maasai.jpg}\hfill\mbox{}

\hfill\includegraphics[width=0.48\textwidth]{my_images/inuit.jpg}\hfill\mbox{}
\uncover<+->\mbox{}\end{frame}

\begin{frame}


\frametitle{{Human heterogeneity}}

\hfill\includegraphics[width=1\textwidth]{my_images/Lesotho_children.jpg}\hfill\mbox{}
\uncover<+->\mbox{}\end{frame}

\begin{frame}


\frametitle{{Human heterogeneity}}

\hfill\includegraphics[height=0.42\textheight]{my_images/dump_workers.jpg}\hfill\mbox{}

\hfill\includegraphics[height=0.48\textheight]{my_images/bento.jpg}\hfill\mbox{}
\uncover<+->\mbox{}\end{frame}


 \section{Homogeneous disease models}
\begin{frame}


\frametitle{Expanding our models}\begin{itemize}[<+-| alert@+>]

\item \textbf{Homogeneous} models assume everyone has the same:\begin{itemize}[<+-| alert@+>]
\item disease characteristics (e.g. susceptibility, tendency to
		transmit)
\item mixing rate
\item probability of mixing with each person\end{itemize}

\item \textbf{Heterogeneous} models allow people to be different\end{itemize}
\uncover<+->\mbox{}\end{frame}

\begin{frame}


\frametitle{The basic reproductive number}\begin{itemize}[<+-| alert@+>]

\item $\Ro$ is the number of people who would be infected by an
	infectious individual \emph{in a fully susceptible population.}

\item $\Ro = \beta/\gamma = \beta D = (cp) D$\begin{itemize}[<+-| alert@+>]
\item $c$: Contact Rate
\item $p$: Probability of transmission (infectivity)
\item $D$: Average duration of infection\end{itemize}

\item A disease can invade a population if and only if $\Ro>1$.\end{itemize}
\uncover<+->\mbox{}\end{frame}

\begin{frame}


\frametitle{Equilibrium}

\hfill\includegraphics[height=0.9\textheight,angle=270]{SIR_model_family/comp.pdf}\hfill\mbox{}\begin{itemize}[<+-| alert@+>]

\item Equilibrium is worth knowing even if the disease doesn't reach equilibrium

\item System will move around the equilibrium\end{itemize}
\uncover<+->\mbox{}\end{frame}

\begin{frame}


\frametitle{Equilibrium analysis }\begin{itemize}[<+-| alert@+>]

\item $\Reffec$ is the number of people who would be infected by an
	infectious individual \emph{in a general population.}

\item $\ds\Reffec =  \Ro  \frac SN = pcD \frac SN$

\item At equilibrium:
		$\ds \Reffec = \Ro  \frac SN = 1. $

\item Thus:
		$\ds \frac SN = 1/R_0.  $

\item Proportion `affected' is $V = 1-S/N = 1-1/R_0$.\end{itemize}
\uncover<+->\mbox{}\end{frame}

\begin{frame}


\frametitle{Proportion affected}\begin{itemize}[<+-| alert@+>]

\item Proportion `affected' is $V = 1-S/N = 1-1/R_0$.\begin{itemize}[<+-| alert@+>]
\item * \uncover<+->{\color{blue} The same formula as the critical vaccination proportion!}
\item * \uncover<+->{\color{blue} If this proportion is made unavailable, the disease cannot spread}
\item * \uncover<+->{\color{blue} At least, in the homogeneous case}\end{itemize}\end{itemize}
\uncover<+->\mbox{}\end{frame}

\begin{frame}


\frametitle{Homogeneous endemic curve}

\begin{columns}[c] \column{0.54\textwidth} \small

\hfill\includegraphics[width=1\textwidth]{Endemic_curves/propCurves.Rout-0.pdf}\hfill\mbox{}

\column{0.4\textwidth} \small\begin{itemize}[<+-| alert@+>]

\item Threshold value

\item Sharp response to changes in factors underlying transmission

\item Works -- sometimes

\item Sometimes predicts unrealistic sensitivity\end{itemize}

\end{columns}
\uncover<+->\mbox{}\end{frame}

\begin{frame}


\frametitle{{Disease dynamics}}

\hfill\includegraphics[height=0.9\textheight]{SIR_simulations/hiv_plot.Rout-0.pdf}\hfill\mbox{}
\uncover<+->\mbox{}\end{frame}

\begin{frame}


\frametitle{{Homogeneous assumptions}}

\hfill\includegraphics[height=0.9\textheight]{SIR_simulations/hiv_plot.Rout-1.pdf}\hfill\mbox{}
\uncover<+->\mbox{}\end{frame}

\begin{frame}


\frametitle{Homogeneous assumptions}

\hfill\includegraphics[height=0.9\textheight]{SIR_simulations/hiv_plot.Rout-2.pdf}\hfill\mbox{}
\uncover<+->\mbox{}\end{frame}

\begin{frame}


\frametitle{Homogeneous assumptions}

\hfill\includegraphics[height=0.9\textheight]{SIR_simulations/hiv_plot.Rout-3.pdf}\hfill\mbox{}
\uncover<+->\mbox{}\end{frame}

\begin{frame}


\frametitle{Homogeneous assumptions}

\hfill\includegraphics[height=0.9\textheight]{SIR_simulations/hiv_plot.Rout-4.pdf}\hfill\mbox{}
\uncover<+->\mbox{}\end{frame}

\begin{frame}


\frametitle{Homogeneous dynamics}\begin{itemize}[<+-| alert@+>]

\item For many diseases, homogeneous models tend to predict:\begin{itemize}[<+-| alert@+>]
\item Too high of an equilibrium, when matching growth rate
\item Too low of a growth rate, when matching equilibrium\end{itemize}\end{itemize}
\uncover<+->\mbox{}\end{frame}


 \section{The importance of heterogeneity}
\begin{frame}


\frametitle{Beyond homogeneity}\begin{itemize}[<+-| alert@+>]

\item Flavors of heterogeneity\begin{itemize}[<+-| alert@+>]
\item among hosts
\item spatial
\item demographic (discreteness of indviduals)
\item temporal
\item others\end{itemize}\end{itemize}
\uncover<+->\mbox{}\end{frame}

\begin{frame}


\frametitle{Heterogeneity in TB}

\begin{columns}[c] \column{0.54\textwidth} \small

\hfill\includegraphics[width=1\textwidth]{Disease_data/tbincome.Rout.pdf}\hfill\mbox{}

\column{0.4\textwidth} \small\begin{itemize}[<+-| alert@+>]

\item {\bf Progression}: Nutrition, stress

\item {\bf Contact}: Overcrowding, poor ventilation

\item {\bf Cure}: Access to medical care\end{itemize}

\end{columns}
\uncover<+->\mbox{}\end{frame}

\begin{frame}


\frametitle{Heterogeneity in other diseases}\begin{itemize}[<+-| alert@+>]

\item {\bf STDs}: Sexual mixing patterns, access to medical care

\item {\bf Influenza}: Crowding, nutrition

\item {\bf Malaria}:  Attractiveness to biting insects, geographical
	location, immune status

\item {\bf Every disease!}\end{itemize}
\uncover<+->\mbox{}\end{frame}

\begin{frame}


\frametitle{Large-scale heterogeneity}

\hfill\includegraphics[height=0.64\textheight]{my_images/schisto.png}\hfill\mbox{}\begin{itemize}[<+-| alert@+>]

\item For schistosomiasis, the worldwide average $\Ro<1$

\item Disease persists because of specific populations with $\Ro>1$.

\item This effect operates at many scales.\end{itemize}
\uncover<+->\mbox{}\end{frame}


\section{Take break} 
\begin{frame}
 \frametitle{Take break}

\hfill\includegraphics[height=0.9\textheight]{webpix/teatime.jpg}\hfill\mbox{}
\uncover<+->\mbox{}\end{frame}


 \section{Effects of heterogeneity}
\begin{frame}


\frametitle{Heterogeneity among hosts}\begin{itemize}[<+-| alert@+>]

\item Differences among people are pervasive, large and often correlated

\item We often consider transmission probability as the product of two
	components:\begin{itemize}[<+-| alert@+>]
\item The "infector" has tendency to infect $\tau$
\item The "recipient" has susceptibility $\sigma$\end{itemize}

\item Then $\Ro = p c D = (\sigma \tau) c D$,

\item Why do we assume this is multiplicative?\begin{itemize}[<+-| alert@+>]
\item * \uncover<+->{\color{blue} Convenience, question this assumption}\end{itemize}\end{itemize}
\uncover<+->\mbox{}\end{frame}

\begin{frame}


\frametitle{Equilibrium calculations}\begin{itemize}[<+-| alert@+>]

\item Assume $p = \sigma\tau$ has a susceptibility component and a
	transmission component:\begin{itemize}[<+-| alert@+>]
\item $\Ro = \sigma\tau c D$
\item $\Reffec = \sigma\tau c D S/N$
\item Equilibrium $S/N = 1/\Ro$
\item Proportion affected: $1 - 1/\Ro$\end{itemize}\end{itemize}
\uncover<+->\mbox{}\end{frame}

\begin{frame}


\frametitle{Equilibrium calculations with heterogeneity}\begin{itemize}[<+-| alert@+>]

\item $\tau D$ applies to infectious individuals $\rightarrow \tau_I
	D_I$

\item $\sigma$ applies to susceptible individuals $\rightarrow
	\sigma_S$

\item $c$ is complicated $\rightarrow c_S c_I/\bar c$\end{itemize}
\uncover<+->\mbox{}\end{frame}

\begin{frame}


\frametitle{Equilibrium calculations with heterogeneity}\begin{itemize}[<+-| alert@+>]

\item $\Ro = \sigma_S\tau_I c_x D_I$
	measured during \emph{invasion}

\item $\Reffec = \sigma_S\tau_I c_x D_I S/N$
	measured at \emph{equilibrium}

\item Equilibrium $S/N \neq 1/\Ro$\end{itemize}
\uncover<+->\mbox{}\end{frame}

\begin{frame}


\frametitle{How does $\Rx$ change?}\begin{itemize}[<+-| alert@+>]

\item Imagine a disease spread by people who differ only in their
	effective mixing rates

\item If the disease has just started spreading in the population, how do
	$c_S$ and $c_I$ compare to $\bar c$?\begin{itemize}[<+-| alert@+>]
\item $c_S \approx \bar c$; $c_I > \bar c$.\end{itemize}

\item If the disease is very widespread in the population?\begin{itemize}[<+-| alert@+>]
\item $c_S < \bar c$; $c_I \to \bar c$.\end{itemize}\end{itemize}
\uncover<+->\mbox{}\end{frame}

\begin{frame}


\frametitle{Simulated population}

\hfill\includegraphics[width=0.8\textwidth]{SIR_model_family/simpson.Rout-0.pdf}\hfill\mbox{}
\uncover<+->\mbox{}\end{frame}

\begin{frame}


\frametitle{Early (5\% infection)}

\hfill\includegraphics[width=0.8\textwidth]{SIR_model_family/simpson.Rout-1.pdf}\hfill\mbox{}
\uncover<+->\mbox{}\end{frame}

\begin{frame}


\frametitle{Mid (20\% infection)}

\hfill\includegraphics[width=0.8\textwidth]{SIR_model_family/simpson.Rout-2.pdf}\hfill\mbox{}
\uncover<+->\mbox{}\end{frame}

\begin{frame}


\frametitle{Mid (50\% infection)}

\hfill\includegraphics[width=0.8\textwidth]{SIR_model_family/simpson.Rout-3.pdf}\hfill\mbox{}
\uncover<+->\mbox{}\end{frame}

\begin{frame}


\frametitle{{\sl Simulated population}}

\hfill\includegraphics[width=0.8\textwidth]{SIR_model_family/simpson.Rout-0.pdf}\hfill\mbox{}
\uncover<+->\mbox{}\end{frame}

\begin{frame}


\frametitle{Simpson's paradox}

\begin{minipage}{\dfwidth} \centering 
\includegraphics[height=0.8\dfheight]{my_images/farmers.jpg} \end{minipage} \hspace{\dfsep} \begin{minipage}{\dfwidth} \centering 
\includegraphics[height=0.8\dfheight]{my_images/senate.jpg} \end{minipage}\begin{itemize}[<+-| alert@+>]

\item What happens when a peanut farmer is elected to the US Senate?

\item The average IQ goes up in both places!\end{itemize}
\uncover<+->\mbox{}\end{frame}

\begin{frame}


\frametitle{The basic reproductive number}\begin{itemize}[<+-| alert@+>]

\item When the disease invades:\begin{itemize}[<+-| alert@+>]
\item The susceptible population $\approx$ the general population
\item The infectious population is likely to have higher values of $c$,
		$D$ and/or $\tau$\end{itemize}

\item $\Ro$ is typically greater than you would expect from a homogeneous
	model\end{itemize}
\uncover<+->\mbox{}\end{frame}

\begin{frame}


\frametitle{Equilibrium analysis}\begin{itemize}[<+-| alert@+>]

\item As disease prevalence goes up:\begin{itemize}[<+-| alert@+>]
\item Susceptible pool is the most resistant, or least exposed group
\item Infectious pool looks more like the general population.\end{itemize}

\item $\rightarrow$ lower proportion affected \emph{for a given value of
	\Ro.}\end{itemize}
\uncover<+->\mbox{}\end{frame}

\begin{frame}


\frametitle{{\sl Homogeneous endemic curve}}

\hfill\includegraphics[height=0.9\textheight]{Endemic_curves/propCurves.Rout-0.pdf}\hfill\mbox{}
\uncover<+->\mbox{}\end{frame}

\begin{frame}


\frametitle{Heterogeneous endemic curves}

\hfill\includegraphics[height=0.9\textheight]{Endemic_curves/propCurves.Rout-2.pdf}\hfill\mbox{}
\uncover<+->\mbox{}\end{frame}

\begin{frame}


\frametitle{Heterogeneous endemic curves}

\hfill\includegraphics[height=0.9\textheight]{Endemic_curves/propCurves.Rout-1.pdf}\hfill\mbox{}
\uncover<+->\mbox{}\end{frame}

\begin{frame}


\frametitle{Heterogeneity and disease}

\begin{columns}[c] \column{0.54\textwidth} \small\begin{itemize}[<+-| alert@+>]

\item Heterogeneity has a double-edged effect\begin{itemize}[<+-| alert@+>]
\item Effects of disease are \emph{lower} for a given value of \Ro.
\item But \Ro\ is \emph{higher} for given mean values of factors
		underlying transmission\end{itemize}\end{itemize}

\column{0.4\textwidth} \small

\hfill\includegraphics[width=1\textwidth]{Endemic_curves/propCurves.Rout-2.pdf}\hfill\mbox{}

\end{columns}
\uncover<+->\mbox{}\end{frame}

\begin{frame}


\frametitle{Heterogeneous endemic curves}

\begin{columns}[c] \column{0.54\textwidth} \small\begin{itemize}[<+-| alert@+>]

\item Heterogeneity makes the endemic curve flatter

\item Disease levels are more resistant to change\end{itemize}

\column{0.4\textwidth} \small

\hfill\includegraphics[width=1\textwidth]{Endemic_curves/propCurves.Rout-1.pdf}\hfill\mbox{}

\end{columns}
\uncover<+->\mbox{}\end{frame}

\begin{frame}


\frametitle{How diseases reach equilibrium }\begin{itemize}[<+-| alert@+>]

\item Diseases that invade have high values of \Ro

\item \Reffec\ must be 1 at equilibrium\begin{itemize}[<+-| alert@+>]
\item Potentially infectious contacts are wasted\begin{itemize}[<+-| alert@+>]
\item Many potential contacts are not susceptible (affected by
			disease)
\item Those not affected less susceptible than average\end{itemize}
\item Infectious pool less infectious\end{itemize}\end{itemize}
\uncover<+->\mbox{}\end{frame}

\begin{frame}


\frametitle{Spatial and network models}\begin{itemize}[<+-| alert@+>]

\item Individual-level, or spatial, heterogeneity also usually
	increases wasted contacts

\item Infectious people meet:\begin{itemize}[<+-| alert@+>]
\item people with similar social backgrounds
\item people with similar behaviours
\item people who are nearby geographically or in the contact network\end{itemize}

\item More wasted contacts further flatten the endemic curve\end{itemize}
\uncover<+->\mbox{}\end{frame}


 \section{Modeling approaches}
\begin{frame}


\frametitle{Phenomenological}

\hfill\includegraphics[height=0.9\textheight]{SIR_model_family/phen.Rout.pdf}\hfill\mbox{}
\uncover<+->\mbox{}\end{frame}

\begin{frame}


\begin{columns}[c] \column{0.54\textwidth} \small

\hfill\includegraphics[height=0.9\textheight]{SIR_model_family/phen.Rout.pdf}\hfill\mbox{}

\column{0.4\textwidth} \small\begin{itemize}[<+-| alert@+>]

\item Simply \emph{make} $\beta$ go down with prevalence, $\beta = B \times$:\begin{itemize}[<+-| alert@+>]
\item $e^{-\alpha P}$
\item $(1-P)^s$
\item $(1-P/s)^{\alpha s}$\end{itemize}\end{itemize}

\end{columns}
\uncover<+->\mbox{}\end{frame}

\begin{frame}


\frametitle{Multi-group models}\begin{itemize}[<+-| alert@+>]

\item Divide the population into groups.\begin{itemize}[<+-| alert@+>]
\item cities and villages
\item rich and poor
\item high and low sexual activity
\item age, gender
\item ...\end{itemize}

\item Even if details are not correct, heterogeneity will emerge and move model in
	the right direction\end{itemize}
\uncover<+->\mbox{}\end{frame}

\begin{frame}


\frametitle{Individual-based models}

\hfill\includegraphics[width=0.8\textwidth]{my_images/network.png}\hfill\mbox{}\begin{itemize}[<+-| alert@+>]

\item Allow many possibilities:\begin{itemize}[<+-| alert@+>]
\item vary individual characteristics
\item add a network of interactions
\item let the network change\end{itemize}

\item Individual-based approaches require stochastic models\end{itemize}
\uncover<+->\mbox{}\end{frame}

\begin{frame}


\frametitle{{Summary}}

\hfill\includegraphics[width=0.4\textwidth]{my_images/crowd.jpg}\hfill\mbox{}

\begin{minipage}{\dfwidth} \centering 
\includegraphics[height=0.8\dfheight]{Endemic_curves/propCurves.Rout-1.pdf} \end{minipage} \hspace{\dfsep} \begin{minipage}{\dfwidth} \centering 
\includegraphics[height=0.8\dfheight]{Endemic_curves/propCurves.Rout-2.pdf} \end{minipage}
\uncover<+->\mbox{}\end{frame}

\begin{frame}


\frametitle{Summary}

\hfill\includegraphics[width=0.32\textwidth]{my_images/crowd.jpg}\hfill\mbox{}\begin{itemize}[<+-| alert@+>]

\item People are heterogeneous in many ways\begin{itemize}[<+-| alert@+>]
\item \ldots and on many scales\end{itemize}

\item Simple models give us important qualitative insights\begin{itemize}[<+-| alert@+>]
\item Diseases in heterogeneous populations are likely to be more robust to
		change than expected from homogeneous models\end{itemize}

\item More complicated models will often be necessary\begin{itemize}[<+-| alert@+>]
\item And it may be helpful to build complexity gradually\end{itemize}\end{itemize}
\uncover<+->\mbox{}\end{frame}

{\topBar%%\sponsorBar

\begin{frame}

\begin{centering}

\vspace{6ex}\includegraphics[height=6ex]{LatexTemplates/rights.png}

\vfill

{
	\tiny\linespread{12}
	This presentation is made available through a Creative Commons Attribution-Noncommercial license.
	Details of the license and permitted uses are available at \url{http://creativecommons.org/licenses/by-nc/3.0/} \par
}

\smallskip
\includegraphics[height=3ex]{LatexTemplates/attrib.png}
\hspace{3em}
\includegraphics[height=3ex]{LatexTemplates/noncom.png}

{\tiny\copyright\ \years, International Clinics on Infectious Disease Dynamics and Data}

\vfill

{\footnotesize
	Title: \inserttitle \\
	Attribution: \insertauthor, \insertdate 
	\par
}

\bigskip

{
	\scriptsize Source URL: \url{\figshare} \\
	For further information please contact \url{admin@ici3d.org}.
	\par
}

\vfill

\end{centering}

\end{frame}

}


\end{document}

